% !TeX root = ../../Book4.tex
% 纲要标题：### 16.1 全纯丛的等价描述
% 纲要位置：几何学计划纲要.md，第 3845 行。
% 正文状态：未撰写；本文件目前是框架。

% 纲要内容开始
% 
% * 全纯过渡函数与局部截面
% * 丛值 ∂̄ 算子和平方为零
% * 可积 ∂̄ 算子的局部平凡化
% 
%   固定光滑复向量丛和光滑 (0,1) 联络，按 Koszul–Malgrange 的局部定理处理 ∂̄²=0；本节安排逐变量求解与规范变换的证明模块，不把形式平方为零直接当作已有全纯标架。
% 
% * 丛的扩张及上同调类
% 
% 纲要内容结束

\section{全纯丛的等价描述}
\label{g4:sec:1601}

待撰写

\subsection{全纯过渡函数与局部截面}
\label{g4:subsec:1601:01}
待撰写

\subsection{\texorpdfstring{丛值 \(\bar{\partial}\) 算子和平方为零}{丛值 \={∂} 算子和平方为零}}
\label{g4:subsec:1601:02}
待撰写

\subsection{\texorpdfstring{可积 \(\bar{\partial}\) 算子的局部平凡化}{可积 \={∂} 算子的局部平凡化}}
\label{g4:subsec:1601:03}
待撰写

\subsection{丛的扩张及上同调类}
\label{g4:subsec:1601:04}
待撰写
% 纲要条目排为编号小节标题，并进入目录与书签；说明段仍保留于纲要注释中。
